\chapter{SEO --- Search Engine Optimisation}
\label{ch:seo}

\section{Applicability}

\textbf{Classic web SEO:} NOT APPLICABLE. The project publishes no website, no
landing page, no HTML documentation site and no marketing URL. There is no
\texttt{index.html}, no \texttt{robots.txt}, no \texttt{sitemap.xml} and no
rendered content for a search engine to index. The only HTTP surface the project
owns is the recipe registry, which serves machine-readable JSON.

\textbf{Discoverability SEO:} APPLICABLE. The product is discovered through three
\textbf{listing} channels that function as its search surface, and these are
where optimisation is meaningful:

\begin{enumerate}
  \item the npm registry (for \texttt{@expo-snips/core} and
        \texttt{@expo-snips/cli});
  \item the VS Code Marketplace listing (for the \texttt{expo-snippets}
        extension);
  \item GitHub repository search and general web search of repository content.
\end{enumerate}

This chapter therefore analyses the metadata that governs discovery in those
channels.

\section{In-repository SEO signals}

\begin{longtable}{p{3.6cm}p{3.2cm}p{7.4cm}}
\caption{SEO-relevant metadata found in the repository}\\
\toprule
\textbf{Artefact} & \textbf{Signal} & \textbf{Content} \\
\midrule

Extension \texttt{package.json} & \texttt{displayName} & ``Expo Snippets''. \\
Extension \texttt{package.json} & \texttt{description} & ``Generate complete React
Native screens with a single abbreviation.'' \\
Extension \texttt{package.json} & \texttt{keywords} & Over 40 keywords covering
\texttt{react-native}, \texttt{expo}, \texttt{screen}, \texttt{template},
\texttt{codegen}, \texttt{expo-router}, \texttt{nativewind}, \texttt{shadcn},
\texttt{vscode}, \texttt{expo-snips} and more. \\
Extension \texttt{package.json} & \texttt{categories} & \texttt{Snippets},
\texttt{Other}. \\
Extension \texttt{package.json} & \texttt{repository.url} & Links to
\texttt{github.com/unkown812/expo-snips}. \\
Extension \texttt{package.json} & \texttt{icon} & \texttt{assets/logo.png}. \\
Extension README & Rich Markdown & Headings, tables, feature lists, quick-start
instructions, badges. \\
CLI \texttt{package.json} & \texttt{description} & ``Beautiful React Native
screens, one command away.'' \\
Core \texttt{package.json} & \texttt{description} & ``Shared types and generation
engine for Expo Snips screens.'' \\
Repository & Topic titles & Recipe names and descriptions provide natural-language
keywords (\texttt{auth/login}, \texttt{onboarding/welcome}, \dots). \\
\bottomrule
\end{longtable}

\section{Strengths}

\begin{itemize}
  \item The extension keyword list is unusually broad and covers the phrases a
        developer would actually search (\texttt{react-native screen generator},
        \texttt{expo template}, \texttt{codegen}, framework names). Keyword
        coverage is the strongest discovery asset the project has.
  \item The extension README is substantial and structured, which is favourable
        both to the Marketplace listing and to GitHub search.
  \item The \texttt{repository} field and README links create a coherent
        cross-channel trail from listing to source.
\end{itemize}

\textbf{Status:} IMPLEMENTED.

\section{Gaps}

\begin{longtable}{p{1.5cm}p{4.6cm}p{3.4cm}p{4.7cm}}
\caption{SEO findings}\\
\toprule
\textbf{ID} & \textbf{Finding} & \textbf{Evidence} & \textbf{Recommendation} \\
\midrule

SEO-01 & No website or landing page exists, so every search result resolves to a
repository or a listing rather than a product page. & Absence of any HTML in the
repo. & If a product identity is intended, publish a small static site with
descriptive copy. \\
SEO-02 & No root README. The top-level repository view shows no descriptive text,
which weakens GitHub search relevance and first impressions. & Absence of
\texttt{README.md} at the repository root. & Add a root README that describes the
monorepo and links the packages. \\
SEO-03 & npm packages carry only the \texttt{description} field as prose; there is
no \texttt{README} referenced from the npm listing and no keywords field on the
npm packages. & \texttt{packages/*/package.json}. & Add \texttt{keywords} and
ensure a README is published with each npm package. \\
SEO-04 & Listing metadata is inconsistent: the extension's repository/author
(\texttt{unkown812}) differs from the owner host implied by some README links
(\texttt{github.com/expo-snips/expo-snips}). & \texttt{package.json} vs README
links. & Standardise the canonical repository URL everywhere. \\
SEO-05 & The registry serves JSON only; nothing about the recipes is
human-browsable or indexable. & \texttt{scripts/serve.mjs}, \texttt{registry/}. &
Optionally publish a human-readable catalog page generated from the index. \\
\bottomrule
\end{longtable}

\section{Summary}

\begin{table}[H]
\centering
\caption{SEO posture}
\begin{tabular}{p{4.6cm}p{9.6cm}}
\toprule
\textbf{Dimension} & \textbf{Assessment} \\
\midrule
Website SEO & NOT APPLICABLE (no website). \\
Listing keyword coverage & IMPLEMENTED and strong (extension keywords). \\
npm discoverability & PARTIAL (descriptions only). \\
Repository discoverability & PARTIAL (no root README). \\
Cross-channel consistency & Needs work (repository URL mismatch). \\
\bottomrule
\end{tabular}
\end{table}

\textbf{Status:} NOT APPLICABLE (web SEO) / PARTIAL (listing SEO) / RECOMMENDED
(root README, canonical URLs).
